\documentclass[12pt, a4paper]{article}
\usepackage[UTF8]{ctex}
\usepackage{amsmath, amssymb}
\usepackage{geometry}
\usepackage{hyperref}
\usepackage{booktabs}
\usepackage{enumitem}
\usepackage{tcolorbox}
\usepackage{fancyhdr}
\usepackage{titlesec}
\usepackage{xcolor}
\usepackage{tabularx}
\usepackage{longtable}

\geometry{left=2.5cm, right=2.5cm, top=2.8cm, bottom=2.8cm}

\definecolor{mainblue}{RGB}{0,82,155}
\definecolor{lightblue}{RGB}{230,240,250}
\definecolor{lightgreen}{RGB}{228,245,228}
\definecolor{lightyellow}{RGB}{255,250,225}
\definecolor{lightred}{RGB}{255,232,232}
\definecolor{lightpurple}{RGB}{242,235,255}

\newtcolorbox{conceptbox}[1][]{%
  colback=lightblue, colframe=mainblue,
  title={\textbf{概念：#1}}, fonttitle=\normalsize,
  boxrule=0.8pt, arc=3pt
}
\newtcolorbox{warningbox}[1][]{%
  colback=lightred, colframe=red!70!black,
  title={\textbf{注意：#1}}, fonttitle=\normalsize,
  boxrule=0.8pt, arc=3pt
}
\newtcolorbox{tipbox}[1][]{%
  colback=lightgreen, colframe=green!50!black,
  title={\textbf{提示：#1}}, fonttitle=\normalsize,
  boxrule=0.8pt, arc=3pt
}
\newtcolorbox{examplebox}[1][]{%
  colback=lightyellow, colframe=orange!70!black,
  title={\textbf{例子：#1}}, fonttitle=\normalsize,
  boxrule=0.8pt, arc=3pt
}
\newtcolorbox{purplebox}[1][]{%
  colback=lightpurple, colframe=purple!60!black,
  title={\textbf{论文原话解读：#1}}, fonttitle=\normalsize,
  boxrule=0.8pt, arc=3pt
}

\pagestyle{fancy}
\fancyhf{}
\fancyhead[L]{\small\textcolor{gray}{碳达峰论文知识点全景解析}}
\fancyhead[R]{\small\textcolor{gray}{从零开始的统计建模入门}}
\fancyfoot[C]{\thepage}
\renewcommand{\headrulewidth}{0.4pt}

\titleformat{\section}{\Large\bfseries\color{mainblue}}{}{0em}{}[\titlerule]
\titleformat{\subsection}{\large\bfseries\color{mainblue!80!black}}{}{0em}{}
\titleformat{\subsubsection}{\normalsize\bfseries\color{mainblue!60!black}}{}{0em}{}

\title{\vspace{-1cm}\textbf{\Huge 碳达峰路径论文知识点全景解析} \\[0.5cm]
\large 基于北京理工大学统计建模大赛参赛作品的完整学习文档 \\[0.3cm]
\normalsize ——从零开始，逐个击破论文中的每一个概念}
\author{}
\date{2026 年}

\begin{document}
\maketitle
\begin{center}
\textcolor{gray}{\small 本文档面向零基础大一新生，将论文中涉及的所有概念逐一拆解讲解。\\
阅读前不需要任何统计学、经济学或计算机基础。}
\end{center}
\newpage
\tableofcontents
\newpage

%% ============================================================
%%  PART I
%% ============================================================
\part{背景与数据篇}

%% ---------- 1 ----------
\section{碳达峰与碳中和：最基本的背景}

\subsection{什么是碳排放？}

人类在燃烧煤炭、石油、天然气等化石燃料时，会向大气中释放二氧化碳（CO$_2$）。此外，某些工业生产过程（如水泥制造）也会排放CO$_2$。这些排放到大气中的CO$_2$总量，就是\textbf{碳排放量}，通常以"百万吨"（Mt，即 Million tonnes）为单位计量。

\begin{examplebox}[生活中的类比]
想象你家里烧煤取暖。每烧一吨煤，就会向空气中释放一定量的CO$_2$。把全中国所有工厂、电厂、汽车、居民……排放的CO$_2$加在一起，就是中国的\textbf{碳排放总量}。
\end{examplebox}

\subsection{什么是"碳达峰"？}

\begin{conceptbox}[碳达峰]
碳排放量不再增长，达到最高点后开始下降。这个最高点就叫做\textbf{"峰值"}，达到峰值的那一年就是\textbf{"达峰年"}。
\end{conceptbox}

用一张简单的图来理解：
\begin{center}
\textit{碳排放量} \\
 $\nearrow$ $\nearrow$ $\rightarrow$ $\searrow$ $\searrow$ \\
\hspace{1cm} \textbf{峰值} $\uparrow$ \end{center}
中国承诺在\textbf{2030 年前}实现碳达峰。

\subsection{什么是"碳中和"？}

\begin{conceptbox}[碳中和]
在一定时期内，排出的CO$_2$与吸收的CO$_2$相互抵消，净排放为零。吸收可以通过植树造林、碳捕集技术等方式实现。
\end{conceptbox}

中国承诺在\textbf{2060 年前}实现碳中和。碳达峰是碳中和的前提——你得先让排放不再增长，才谈得上最终降到零。

\subsection{为什么要研究碳达峰路径？}

\begin{itemize}[leftmargin=2em]
\item \textbf{什么时候达峰？}——太晚达峰意味着留给碳中和的时间不够。
\item \textbf{峰值有多高？}——峰值越高，后面下降越困难，成本越大。
\item \textbf{是什么因素在驱动排放变化？}——搞清楚"病因"才能"对症下药"。
\item \textbf{不同政策下未来会怎样？}——帮政府做决策参考。
\end{itemize}

这就是这篇论文要回答的核心问题。

%% ---------- 2 ----------
\section{数据从哪来？怎么整理？}

\subsection{CEADs 数据库}

\begin{conceptbox}[CEADs]
全称 \textit{Carbon Emission Accounts \& Datasets}（碳核算与数据集），是中国学者构建的权威碳排放数据库。它提供中国各省、各年份的碳排放数据，数据来源于化石燃料燃烧和工业过程排放的估算。
\end{conceptbox}

\begin{tipbox}[论文怎么用 CEADs 的？]
CEADs 给出的是\textbf{省级}数据（比如北京多少、河北多少……）。论文需要的是\textbf{国家级}数据，所以做法很简单——把31个省份同年份的排放量\textbf{加起来}：
\[
C_t = \sum_{i=1}^{N} C_{i,t}
\]
其中 $C_{i,t}$ 是第 $i$ 个省第 $t$ 年的碳排放量。这就是论文公式(1)。
\end{tipbox}

\subsection{国家统计局数据}

除了碳排放量（因变量），论文还需要很多"解释变量"来解释为什么排放会变化。这些数据来自\textbf{国家统计局}年度数据库：
\begin{itemize}[leftmargin=2em]
\item 人口总数
\item GDP（国内生产总值）
\item 城镇化率
\item 第二产业占比
\item 能源消费总量
\item 煤炭消费占比
\item 专利数据（代理技术创新）
\end{itemize}

\subsection{什么是"口径统一"？为什么重要？}

\begin{warningbox}[口径统一——大一新生最容易忽视的问题]
假设你想比较2010年和2020年的GDP。但2010年的GDP是用"2000年不变价"算的，2020年用的是"2015年不变价"，两套标准不一样，直接比较就会出错。\textbf{口径统一}就是确保所有年份的数据用同一套标准来度量。
\end{warningbox}

论文提到的三条原则：
\begin{enumerate}[leftmargin=2em]
\item 同一变量只用一个来源，不混用
\item 优先用连续年份完整的来源
\item 必要时做单位换算，并在分析前做一致性检查
\end{enumerate}

\subsection{论文中的变量一览}

\begin{table}[H]
\centering
\small
\begin{tabularx}{\textwidth}{lllX}
\toprule
\textbf{符号} & \textbf{名称} & \textbf{单位} & \textbf{直觉含义} \\
\midrule
 $C_t$ & 碳排放总量 & Mt & 我们要解释和预测的对象 \\
 $P_t$ & 总人口 & 万人 & 人越多，排放越多？ \\
 $A_t$ & 人均GDP & 万元/人 & 越富裕，排放越多？ \\
 $EI_t$ & 能源强度 & 万吨标煤/亿元 & 每单位GDP消耗多少能源 \\
 $CS_t$ & 煤炭占比 & 比例 & 能源中煤炭占多少 \\
 $UR_t$ & 城镇化率 & 比例 & 城镇人口占总人口比 \\
 $IS_t$ & 第二产业占比 & 比例 & 工业占GDP的比重 \\
 $TP_t$ & 技术进步代理 & — & 用专利数等代表技术水平 \\
\bottomrule
\end{tabularx}
\caption{论文核心变量一览}
\end{table}

\subsection{几个变量的计算公式详解}

\subsubsection{人均GDP（论文公式2）}
\[
A_t = \frac{\text{GDP}_t}{P_t}
\]
就是把总GDP除以总人口。这衡量的是"平均每个人创造的经济产出"。

\subsubsection{能源强度（论文公式3）}
\[
EI_t = \frac{EC_t}{\text{GDP}_t}
\]
 $EC_t$ 是能源消费总量。能源强度的含义是：\textbf{每创造1亿元GDP，需要消耗多少万吨标准煤}。这个值\textbf{越小越好}——意味着同样的经济产出只需要更少的能源。

\subsubsection{碳强度（论文公式4）}
\[
CI_t = \frac{C_t}{\text{GDP}_t}
\]
每创造1亿元GDP，排放多少Mt碳。也是越小越好。

\subsubsection{人均碳排放（论文公式5）}
\[
PC_t = \frac{C_t}{P_t}
\]
平均每个人排放多少碳。

\subsection{描述性统计：看数据"长什么样"}

\begin{conceptbox}[描述性统计]
用几个简单的数字来概括一批数据的整体特征，最常见的有：
\begin{itemize}[leftmargin=2em]
\item \textbf{均值}：所有值的平均
\item \textbf{标准差}：数据偏离均值的程度，越大越分散
\item \textbf{最小值/最大值}：数据的范围
\end{itemize}
\end{conceptbox}

\begin{examplebox}[论文中的描述性统计]
碳排放总量在2005--2022年间：均值约30566.50 Mt，最小17516.39 Mt，最大40358.26 Mt。说明排放总量在这个期间\textbf{翻了一倍多}，整体呈显著上升趋势。
\end{examplebox}

\subsection{相关性矩阵热力图（图3）}

\begin{conceptbox}[相关系数]
衡量两个变量之间线性关系的强弱。取值范围 $[-1, 1]$：
\begin{itemize}[leftmargin=2em]
\item 接近 $+1$：一个变大，另一个也变大（正相关）
\item 接近 $-1$：一个变大，另一个变小（负相关）
\item 接近 $0$：没有明显的线性关系
\end{itemize}
\end{conceptbox}

\begin{tipbox}[论文图3告诉我们什么？]
热力图中颜色越深（红或蓝），说明相关性越强。论文发现：GDP、能源消费等\textbf{规模变量}与碳排放正相关（红色）；煤炭占比下降与减排方向一致（负相关）。这为后面选变量提供了依据。
\end{tipbox}

\subsection{对数变换：为什么要把数据取对数？}

论文中大量使用 $\ln$（自然对数），这是一个非常关键的数据预处理步骤。

\begin{conceptbox}[对数变换]
把原始值 $x$ 变成 $\ln(x)$。比如 $\ln(100) \approx 4.61$，$\ln(10000) \approx 9.21$。
\end{conceptbox}

\textbf{为什么要取对数？三个原因：}

\begin{enumerate}[leftmargin=2em]
\item \textbf{把"乘法关系"变成"加法关系"}。比如 $I = P \times A \times T$，取对数后变成 $\ln I = \ln P + \ln A + \ln T$，就可以用线性回归来估计了。
\item \textbf{缓解异方差}。经济数据通常"大的值波动更大，小的值波动更小"。取对数后，波动变得更均匀。
\item \textbf{系数可以直接解释为弹性}（下文详述）。
\end{enumerate}

\begin{warningbox}[什么变量取对数，什么不取？]
论文的策略是：
\begin{itemize}[leftmargin=2em]
\item \textbf{规模变量}（人口、GDP、能源强度等）$\to$ 取对数
\item \textbf{比例变量}（煤炭占比、城镇化率、第二产业占比等）$\to$ \textbf{不取}对数，保持原始比例形式
\end{itemize}
因为比例变量已经在 $[0,1]$ 范围内，取对数反而会扭曲含义。
\end{warningbox}

\subsection{标准化}

\begin{conceptbox}[标准化]
把变量变成均值为0、标准差为1的形式：
\[
x_{\text{new}} = \frac{x - \bar{x}}{s}
\]
其中 $\bar{x}$ 是均值，$s$ 是标准差。
\end{conceptbox}

\begin{tipbox}[论文中标准化用在哪里？]
只在\textbf{机器学习模型}（XGBoost、随机森林）输入前做标准化。因为不同变量的数量级差异很大（GDP是几十万亿，城镇化率是小数），标准化可以让模型训练更稳定。OLS回归没有做标准化，因为对数变换已经缓解了量级问题。
\end{tipbox}

\subsection{缺失值处理与样本边界}

\begin{warningbox}[论文的一个关键区分]
\begin{itemize}[leftmargin=2em]
\item \textbf{2005--2022年}：有真实观测数据，用于\textbf{机制识别}（回归分析）
\item \textbf{2023--2024年}：没有真实碳排放数据（CEADs还没更新到），仅用于预测的\textbf{趋势衔接}，\textbf{不参与}回归估计
\end{itemize}
这个区分非常重要——不能把模型预测出来的值当"真实数据"来用，否则就犯了"用未来解释过去"的错误。
\end{warningbox}

%% ============================================================
%%  PART II
%% ============================================================
\part{统计基础篇}

%% ---------- 3 ----------
\section{OLS回归：一切的基础}

\subsection{什么是回归？}

\begin{conceptbox}[回归分析]
研究一个"结果变量"（因变量 $Y$）是如何被若干"原因变量"（自变量 $X_1, X_2, \ldots$）影响的。简单说就是：\textbf{找关系}。
\end{conceptbox}

\begin{examplebox}[生活中的回归]
假设你想研究"学习时间"对"考试分数"的影响。你收集了100个同学的数据，画散点图，发现大致是一条直线关系。这条直线就是"回归线"，它的方程就是"回归方程"。
\end{examplebox}

\subsection{什么是OLS？}

\begin{conceptbox}[OLS（普通最小二乘法）]
\textit{Ordinary Least Squares}的缩写。它是找回归线的一种方法：找到一条线，使得所有数据点到这条线的\textbf{垂直距离的平方和}最小。
\end{conceptbox}

数学上，对于模型 $Y = \beta_0 + \beta_1 X + \varepsilon$，OLS的目标是：
\[
\min_{\beta_0, \beta_1} \sum_{i=1}^{n} (Y_i - \hat{\beta}_0 - \hat{\beta}_1 X_i)^2
\]
这里的"二乘"就是"平方"的意思。

\subsection{多元回归}

当有多个自变量时：
\[
Y = \beta_0 + \beta_1 X_1 + \beta_2 X_2 + \cdots + \beta_k X_k + \varepsilon
\]

\begin{tipbox}[论文中的STIRPAT就是多元回归]
论文的扩展STIRPAT模型（公式12）有6个自变量，就是多元回归。每个 $\beta$ 代表"在\textbf{控制其他变量不变}的情况下，该变量对碳排放的影响"。
\end{tipbox}

\subsection{弹性}

\begin{conceptbox}[弹性]
一个变量变化1\%时，另一个变量变化百分之几。
\end{conceptbox}

\begin{examplebox}[弹性的直觉]
如果人均GDP的弹性是0.59，意思是：人均GDP增加1\%，碳排放增加0.59\%。注意是\textbf{百分比变化}，不是绝对量变化。
\end{examplebox}

\textbf{关键点：对数回归中系数就是弹性。}

如果模型是 $\ln Y = \beta \ln X + \cdots$，那么：
\[
\beta = \frac{\Delta \ln Y}{\Delta \ln X} \approx \frac{\Delta Y / Y}{\Delta X / X}
\]
这就是弹性的定义！所以取对数的一个好处就是\textbf{系数直接读作弹性}。

\subsection{半弹性}

\begin{conceptbox}[半弹性]
自变量变化\textbf{1个单位}（不是1\%）时，因变量变化百分之几。
\end{conceptbox}

论文中煤炭占比 $CS_t$ 没有取对数，而 $C_t$ 取了对数，所以 $CS_t$ 的系数是半弹性：煤炭占比每增加1个百分点（比如从60\%到61\%），碳排放变化百分之几。

\subsection{异方差}

\begin{conceptbox}[异方差]
回归模型假设误差项的方差是常数（\textit{同方差}）。如果误差的方差随自变量变化而变化，就叫\textit{异方差}。
\end{conceptbox}

\begin{examplebox}[异方差的直觉]
低收入家庭消费的波动小（吃饭穿衣基本固定），高收入家庭消费的波动大（可以奢侈也可以节俭）。用消费做因变量时，就存在异方差。
\end{examplebox}

\begin{tipbox}[对数变换如何缓解异方差？]
取对数把数据的"尺度"压缩了。原来从100到10000差100倍，取对数后从4.6到9.2只差2倍，波动变得更均匀。
\end{tipbox}

\subsection{多重共线性}

\begin{conceptbox}[多重共线性]
多个自变量之间高度相关，导致回归系数估计不稳定（标准误变大）。
\end{conceptbox}

\begin{warningbox}[论文中的巧妙处理]
人口 $P_t$ 和人均GDP $A_t = \text{GDP}/P_t$ 之间天然相关（因为人均GDP就是用人口算出来的）。如果同时把人口和人均GDP放进回归，会造成严重共线性。论文选择\textbf{只用人均GDP，不放人口}，巧妙避开了这个问题。
\end{warningbox}

%% ---------- 4 ----------
\section{统计推断：我们怎么知道结果"靠谱"？}

\subsection{置信区间}

\begin{conceptbox}[95\%置信区间]
如果我们重复抽样100次，每次都算一个置信区间，大约有95个区间会包含真实值。注意：不是说"真实值有95\%的概率落在这个区间里"（真实值是固定的，不是随机的）。
\end{conceptbox}

论文中系数估计的置信区间形式为（公式13）：
\[
\hat{\beta}_j \pm z_{0.975} \cdot \frac{\hat{\sigma}}{\sqrt{[(X^\top X)^{-1}]_{jj}}}
\]
\begin{itemize}[leftmargin=2em]
\item $\hat{\beta}_j$：系数的点估计值
\item $z_{0.975} \approx 1.96$：标准正态分布的97.5\%分位数
\item 后面的部分：标准误（衡量估计的不确定性）
\end{itemize}

\begin{examplebox}[论文中的实例]
煤炭占比的弹性估计值是3.3791，95\%置信区间是 $[0.0469, 6.7112]$。这个区间\textbf{不包含0}，说明煤炭占比对碳排放的影响在统计上是显著的。

而人均GDP的弹性是0.5882，置信区间是 $[-0.9264, 2.1027]$。这个区间\textbf{包含0}，说明在当前样本下，我们不能确定人均GDP对碳排放到底是正影响还是负影响。
\end{examplebox}

\subsection{显著性检验与p值}

\begin{conceptbox}[p值]
在"该变量对因变量没有影响"（原假设）的前提下，观察到当前数据或更极端数据的概率。p值越小，越有理由拒绝原假设，即认为该变量"确实有影响"。
\end{conceptbox}

\begin{itemize}[leftmargin=2em]
\item $p < 0.05$：在95\%置信水平下显著（论文用的标准）
\item $p < 0.01$：在99\%置信水平下显著
\item $p \geq 0.05$：不显著
\end{itemize}

\begin{tipbox}[置信区间与显著性检验的关系]
\textbf{等价判断}：95\%置信区间不包含0 $\iff$ $p < 0.05$。论文用置信区间来判断显著性，和用p值是一回事。
\end{tipbox}

\subsection{自由度}

\begin{conceptbox}[自由度]
可以"自由变化"的数据个数。在回归中：
\[
\text{自由度} = n - k
\]
其中 $n$ 是样本量，$k$ 是要估计的参数个数（包括截距项）。
\end{conceptbox}

\begin{warningbox}[论文的核心痛点：自由度严重不足！]
论文有18年数据（2005--2022），要估计7个参数（6个自变量 + 1个截距），所以自由度只有 $18 - 7 = 11$。

这\textbf{非常少}。自由度越小，置信区间越宽，越难得到"显著"的结果。论文中很多变量不显著，\textbf{不是}因为它们真的没影响，而是因为数据太少，统计检验"功力不够"。
\end{warningbox}

\begin{examplebox}[类比]
想象你要判断一枚硬币是否均匀。抛10次，可能还看不出规律。抛10000次，就很有把握了。自由度就像"抛硬币的次数"——次数越多，判断越准确。论文只有11个自由度，相当于只抛了11次。
\end{examplebox}

%% ============================================================
%%  PART III
%% ============================================================
\part{机制识别模型篇}

%% ---------- 5 ----------
\section{IPAT恒等式：一切的起点}

\subsection{IPAT是什么？}

\begin{conceptbox}[IPAT恒等式]
环境影响（$I$）= 人口（$P$）$\times$ 富裕度（$A$）$\times$ 技术（$T$）：
\[
I = P \times A \times T
\]
这是一个\textbf{恒等式}（等号一定成立），不是回归模型。它只是把环境影响分解成三个因素的乘积。
\end{conceptbox}

\begin{itemize}[leftmargin=2em]
\item $I$（Impact）：环境影响，这里指碳排放
\item $P$（Population）：人口
\item $A$（Affluence）：富裕度，通常用人均GDP
\item $T$（Technology）：技术水平，通常定义为 $I/(P \times A)$，即单位人口单位GDP的排放
\end{itemize}

\begin{warningbox}[IPAT的局限]
IPAT假设三个因素的\textbf{弹性都是1}（即人口增加1\%，排放一定增加1\%）。这太死板了。现实中弹性可能大于1也可能小于1。IPAT不能做统计检验，只是一个会计分解。
\end{warningbox}

%% ---------- 6 ----------
\section{STIRPAT模型：论文的核心解释模型}

\subsection{从IPAT到STIRPAT}

\begin{conceptbox}[STIRPAT]
\textit{Stochastic Impacts by Regression on Population, Affluence, and Technology}。它把IPAT的固定关系变成\textbf{随机}的、可估计的\textbf{回归}模型。
\end{conceptbox}

具体做法：给每个因素加上指数参数，再加一个随机误差：
\[
I = a \cdot P^b \cdot A^c \cdot T^d \cdot e \quad \text{（论文公式10）}
\]
\begin{itemize}[leftmargin=2em]
\item $a$ 是常数系数
\item $b, c, d$ 是\textbf{弹性参数}——不再固定为1，而是由数据估计出来
\item $e$ 是随机误差
\end{itemize}

两边取对数得到（论文公式11）：
\[
\ln I_t = \ln a + b \ln P_t + c \ln A_t + d \ln T_t + \ln e_t
\]
这就可以用OLS来估计了！$\hat{b}, \hat{c}, \hat{d}$ 就是各个因素的弹性。

\subsection{论文的"扩展"STIRPAT}

\begin{conceptbox}[扩展STIRPAT]
在标准三因素（人口、富裕度、技术）基础上，加入更多与现实相关的变量，使模型更贴合中国实际。
\end{conceptbox}

论文的扩展形式（公式12）：
\[
\ln C_t = \beta_0 + \beta_1 \ln A_t + \beta_2 \ln EI_t + \beta_3 CS_t + \beta_4 UR_t + \beta_5 TP_t + \beta_6 IS_t + \varepsilon_t
\]

\begin{table}[H]
\centering
\small
\begin{tabularx}{\textwidth}{llcl}
\toprule
\textbf{变量} & \textbf{含义} & \textbf{预期方向} & \textbf{为什么？} \\
\midrule
 $\ln A_t$ & 人均GDP & $+$ & 越富裕，消耗越多能源 \\
 $\ln EI_t$ & 能源强度 & $+$ & 单位GDP能耗越高，排放越多 \\
 $CS_t$ & 煤炭占比 & $+$ & 煤是最"脏"的化石能源 \\
 $UR_t$ & 城镇化率 & $\pm$ & 可能增加（基建），也可能减少（集约化） \\
 $TP_t$ & 技术进步 & $-$ & 技术进步提高效率，减少排放 \\
 $IS_t$ & 第二产业占比 & $+$ & 工业是高排放行业 \\
\bottomrule
\end{tabularx}
\caption{各变量的预期影响方向}
\end{table}

\subsection{论文STIRPAT的核心结果解读}

\begin{purplebox}[论文发现：只有煤炭占比显著！]
在95\%置信水平下，\textbf{只有煤炭占比}的弹性显著不为零（弹性约3.3791）。其他变量虽然点估计方向与直觉一致，但置信区间太宽、包含了零，不够显著。

\textbf{怎么理解弹性3.3791？}煤炭占比每增加1个百分点，碳排放增加约3.38\%。这个影响非常大！

\textbf{为什么其他变量不显著？}核心原因是\textbf{自由度太少}（只有11），不是它们没影响。论文明确说了："并不否定其经济机制上的现实作用"。
\end{purplebox}

%% ---------- 7 ----------
\section{EKC环境库兹涅茨曲线}

\subsection{什么是库兹涅茨曲线？}

\begin{conceptbox}[库兹涅茨曲线（原始版）]
经济学家Kuznets发现：随着经济发展，收入不平等先变大后变小，呈倒U型。
\end{conceptbox}

\subsection{什么是环境库兹涅茨曲线（EKC）？}

\begin{conceptbox}[EKC假说]
把这个逻辑搬到环境领域：随着人均收入增长，环境污染先恶化后改善，呈\textbf{倒U型}。
\end{conceptbox}

直觉理解：
\begin{enumerate}[leftmargin=2em]
\item \textbf{穷的时候}：首要目标是发展经济，不太管环境 $\to$ 排放增加
\item \textbf{中等收入时}：工业高速发展 $\to$ 排放达到最高
\item \textbf{富裕之后}：有财力和技术来治理环境 $\to$ 排放下降
\end{enumerate}

\subsection{EKC的数学形式（论文公式6）}

\[
\ln C_t = \alpha_0 + \alpha_1 \ln A_t + \alpha_2 (\ln A_t)^2 + \varepsilon_t
\]

这是一个关于 $\ln A_t$ 的\textbf{二次函数}。当 $\alpha_1 > 0$ 且 $\alpha_2 < 0$ 时，曲线呈倒U型。

\subsection{如何找"拐点"？（论文公式7、8）}

对 $\ln A_t$ 求导，令导数为零：
\[
\frac{\partial \ln C_t}{\partial \ln A_t} = \alpha_1 + 2\alpha_2 \ln A_t = 0
\]
解出拐点处的人均GDP：
\[
A^* = \exp\left(-\frac{\alpha_1}{2\alpha_2}\right)
\]

\begin{examplebox}[拐点的含义]
如果拐点对应人均GDP为5万元，意思是：当人均GDP超过5万元时，经济增长不再带来排放增加，反而排放开始下降。中国是否已经过了这个拐点？论文通过数据来检验。
\end{examplebox}

\begin{tipbox}[论文的定位]
论文明确说EKC检验的作用不是"证明因果关系"，而是判断"经济规模扩张是否已出现边际减排拐点\textbf{迹象}"。这是一个非常谨慎的学术表述。
\end{tipbox}

%% ---------- 8 ----------
\section{中介效应分析}

\subsection{什么是中介效应？}

\begin{conceptbox}[中介效应]
自变量 $X$ 影响因变量 $Y$，但这个影响可能是\textbf{通过}中间变量 $M$ 间接实现的。$M$ 就是"中介变量"。
\end{conceptbox}

\begin{examplebox}[生活例子]
\textbf{锻炼}（$X$）$\to$ \textbf{减轻体重}（$Y$）。但锻炼不是直接让你变瘦的，而是通过 $\to$ \textbf{消耗卡路里}（$M$）$\to$ 减轻体重。这里"消耗卡路里"就是中介变量。
\end{examplebox}

\subsection{中介效应的三步法（论文公式14--16）}

\begin{enumerate}[leftmargin=2em]
\item \textbf{a路径}：$X$ 是否影响 $M$？
\[
M_t = \gamma_0 + \gamma_1 X_t + u_t
\]
如果 $\gamma_1$ 显著，说明 $X$ 确实影响 $M$。

\item \textbf{c路径（总效应）}：$X$ 是否影响 $Y$？（不考虑 $M$）
\[
Y_t = \delta_0 + \delta_1 X_t + v_t
\]
 $\delta_1$ 是总效应。

\item \textbf{c'路径 + b路径}：控制 $M$ 后，$X$ 对 $Y$ 的直接效应，以及 $M$ 对 $Y$ 的效应：
\[
Y_t = \theta_0 + \theta_1 X_t + \theta_2 M_t + w_t
\]
\end{enumerate}

\subsection{效应分解（论文公式17）}

\[
\underbrace{\delta_1}_{\text{总效应}} = \underbrace{\theta_1}_{\text{直接效应}} + \underbrace{\gamma_1 \cdot \theta_2}_{\text{间接效应（中介效应）}}
\]

\begin{itemize}[leftmargin=2em]
\item \textbf{直接效应} $\theta_1$：不通过 $M$ 的那部分影响
\item \textbf{间接效应} $\gamma_1 \cdot \theta_2$：通过 $M$ 传导的那部分
\item 如果 $|\theta_1| < |\delta_1|$ 且间接效应显著 $\to$ \textbf{部分中介}
\item 如果 $\theta_1$ 不显著且间接效应显著 $\to$ \textbf{完全中介}
\end{itemize}

\subsection{论文中的中介分析结果}

\begin{purplebox}[一个非常精彩的发现："符号反转"]
论文设 $X$ = 煤炭占比，$M$ = 技术进步，$Y$ = 碳排放。

\textbf{奇怪的现象}：在不控制GDP时，煤炭占比对碳排放的系数是\textbf{负的}（$c = -3.31$）！这违反直觉——煤炭占比越高，排放应该越多才对。

\textbf{原因}：2005--2022年间，中国GDP和碳排放在涨，但煤炭占比在降。简单回归把这种"时间上的反向趋势"误读成了因果关系。

\textbf{控制GDP后}：煤炭占比的直接效应变成\textbf{正的}（$c' = 2.22$），符合物理直觉！这就是\textbf{伪相关}的经典案例，也说明了多变量控制的重要性。
\end{purplebox}

\begin{tipbox}[中介分析的核心结论]
能源替代（降煤炭占比）的减排逻辑\textbf{主要靠直接效应}（直接压降高碳底座），\textbf{不是}靠间接的技术创新传导。间接效应虽然有，但太小了。所以政策应该直接盯着"降煤"这个目标。
\end{tipbox}

\subsection{Kaya恒等式与LMDI分解（文献综述中提到）}

\begin{conceptbox}[Kaya恒等式]
把碳排放分解为四个因素的乘积：
\[
C = \text{人口} \times \frac{\text{GDP}}{\text{人口}} \times \frac{\text{能源}}{\text{GDP}} \times \frac{\text{碳排放}}{\text{能源}}
\]
即：$C = P \times \text{人均GDP} \times \text{能源强度} \times \text{碳排放系数}$ \end{conceptbox}

\begin{conceptbox}[LMDI分解]
\textit{Logarithmic Mean Divisia Index}。一种把碳排放的\textbf{变化量}分解为各因素\textbf{贡献量}的方法。比如今年排放比去年多了100Mt，LMDI可以告诉你：其中50Mt是因为GDP增长，$-30Mt$是因为能源强度下降，等等。
\end{conceptbox}

\begin{tipbox}[LMDI vs STIRPAT]
\begin{itemize}[leftmargin=2em]
\item LMDI：做\textbf{事后归因}——"过去的变化是怎么来的"
\item STIRPAT：做\textbf{边际效应估计}——"各因素的弹性有多大"
\end{itemize}
论文选择了STIRPAT路线，因为更适合与预测和情景模拟衔接。
\end{tipbox}

%% ============================================================
%%  PART IV
%% ============================================================
\part{时间序列基础篇}

%% ---------- 9 ----------
\section{时间序列的基本概念}

\subsection{什么是时间序列数据？}

\begin{conceptbox}[时间序列]
按时间顺序排列的数据。比如中国2005年的碳排放、2006年的碳排放……直到2022年。每一年只有一个数据点，这就是\textbf{时间序列}（区别于\textbf{截面数据}——同一时间多个个体的数据，和\textbf{面板数据}——多个个体多年的数据）。
\end{conceptbox}

\subsection{平稳性}

\begin{conceptbox}[平稳时间序列]
均值、方差、自相关结构\textbf{不随时间变化}的序列。直观地说，就是"没有趋势、没有周期性波动变化，始终在一个固定水平附近波动"。
\end{conceptbox}

\begin{examplebox}[平稳 vs 不平稳]
\begin{itemize}[leftmargin=2em]
\item \textbf{平稳}：每天抛硬币的结果序列。均值始终是0.5，方差始终不变。
\item \textbf{不平稳}：中国GDP序列。每年都在涨，均值每年都在变。
\end{itemize}
\end{examplebox}

\subsection{单位根与I(1)过程}

\begin{conceptbox}[单位根]
一种导致不平稳的"根源"。如果一个序列有单位根，对它做任何冲击，影响都会\textbf{永久持续}，不会消失。
\end{conceptbox}

\begin{conceptbox}[I(1)过程——"一阶单整"]
原始序列不平稳（有单位根），但做\textbf{一阶差分}后变平稳。一阶差分就是用今年的值减去去年的值：$\Delta C_t = C_t - C_{t-1}$。
\end{conceptbox}

\begin{examplebox}[I(1)的直觉]
中国的碳排放每年都在涨，不平稳。但"每年比上一年多排了多少"，这个变化量大致在一个范围内波动，可能是平稳的。这就叫I(1)。
\end{examplebox}

\subsection{单位根检验}

\begin{conceptbox}[单位根检验]
用统计方法检验一个序列是否有单位根（即是否不平稳）。常见的方法有ADF检验、PP检验等。
\end{conceptbox}

\begin{itemize}[leftmargin=2em]
\item 原假设 $H_0$：有序列单位根（不平稳）
\item 备择假设 $H_1$：无单位根（平稳）
\item 如果 $p < 0.05$，拒绝原假设，认为序列平稳
\end{itemize}

\begin{purplebox}[论文的发现]
论文提到"单位根检验表明序列为I(1)过程"，即碳排放序列不平稳、有一阶单位根。这解释了为什么AR模型（利用差分/滞后信息）比水平拟合的复杂模型更稳健。
\end{purplebox}

\subsection{伪回归}

\begin{warningbox}[伪回归——时间序列分析中的头号陷阱]
当两个\textbf{都不平稳}的序列做回归时，即使它们之间\textbf{完全没有关系}，回归结果也可能表现出很高的$R^2$和"显著"的系数。这就是伪回归。
\end{warningbox}

\begin{examplebox}[经典例子]
有人拿"中国的碳排放"和"美国某条河的水位"做回归，可能发现"显著正相关"。但这纯粹是因为两个序列都在涨（都有趋势），不是真有关系。
\end{examplebox}

\subsection{协整}

\begin{conceptbox}[协整]
两个不平稳的I(1)序列，如果存在某个\textbf{线性组合}是平稳的，就说它们有协整关系。协整意味着两个序列之间存在\textbf{真实的长期均衡关系}，不是伪回归。
\end{conceptbox}

\begin{warningbox}[论文为什么没做协整检验？]
论文明确说明：样本量只有18，传统Johansen协整检验在这种极小样本下"统计功效不足"（即检验能力太弱，分不出真假）。所以论文把STIRPAT的结果定位为"弹性方向参考"而非"严谨因果推断"，把最终结论放在不需要平稳性假设的滚动预测和情景模拟上。这是一个\textbf{非常诚实的学术处理}。
\end{warningbox}

%% ============================================================
%%  PART V
%% ============================================================
\part{预测模型篇}

%% ---------- 10 ----------
\section{滚动一步预测：论文的评估方法}

\subsection{什么是"滚动一步预测"？}

\begin{conceptbox}[滚动一步向前预测]
预测第 $t$ 期时，\textbf{只用}第1到第 $t-1$ 期的数据来训练模型，然后预测第 $t$ 期。接着窗口向前滚动一期，用1到第 $t$ 期的数据预测第 $t+1$ 期。如此反复。
\end{conceptbox}

\begin{examplebox}[与"全样本拟合"的区别]
\textbf{全样本拟合}：用2005--2022年全部数据训练，然后"预测"这些年本身。这相当于考试时提前看了答案，当然分数很高。

\textbf{滚动预测}：用2005--2010年数据预测2011年，再用2005--2011年预测2012年……每次预测都是"真刀真枪"的，不能偷看答案。
\end{examplebox}

\begin{tipbox}[为什么论文选择滚动预测？]
因为论文关心的是\textbf{模型对未来新数据的预测能力}（外推能力），而不是对历史数据的拟合能力。滚动预测严格避免了"未来信息泄露"，是最诚实的评估方式。
\end{tipbox}

%% ---------- 11 ----------
\section{七种预测模型逐一讲解}

\subsection{模型1：朴素模型（Naive）（论文公式18）}

\[
\hat{C}_t = C_{t-1}
\]

\begin{conceptbox}[朴素模型]
最简单的预测方法：明天的排放量 = 今天的排放量。假设排放遵循"随机游走"——没有趋势，没有规律，纯随机波动。
\end{conceptbox}

\begin{tipbox}[朴素模型的作用]
它是一个\textbf{基准}。如果你的复杂模型连朴素模型都打不过，那这个复杂模型就是废的。就像你训练了一个AI模型来做预测，结果还不如"直接猜昨天的值"，那这个AI就没意义了。
\end{tipbox}

\subsection{模型2：二阶自回归模型 AR2（论文公式19）}

\[
\hat{C}_t = \phi_0 + \phi_1 C_{t-1} + \phi_2 C_{t-2}
\]

\begin{conceptbox}[自回归模型（AR模型）]
用变量\textbf{过去的值}来预测\textbf{未来的值}。AR2就是用前两期的值来预测下一期。
\end{conceptbox}

\begin{examplebox}[AR2的直觉]
"过去两年排放分别是38000和39000，按照这个惯性，今年大概是……"AR2就是把这个"惯性"用数学公式表达出来。$\phi_1$ 和 $\phi_2$ 由OLS在训练窗口内估计。
\end{examplebox}

\begin{purplebox}[论文的冠军模型]
AR2在滚动预测中表现\textbf{最好}：RMSE=1323.36，MAPE=3.32\%。一个简单的线性模型打败了所有复杂的机器学习模型！这说明了什么？\textbf{在数据极少的情况下，简单模型往往更可靠。}
\end{purplebox}

\subsection{模型3：二次趋势模型（论文公式20）}

\[
\hat{C}_t = \eta_0 + \eta_1 t + \eta_2 t^2
\]

\begin{conceptbox}[二次趋势模型]
假设碳排放随时间呈\textbf{抛物线}趋势。$\eta_2 > 0$ 表示加速上升，$\eta_2 < 0$ 表示增速放缓或出现拐点。
\end{conceptbox}

\begin{warningbox}[趋势外推的危险]
这个模型假设未来会延续历史的趋势形状。但如果趋势发生"结构性变化"（比如突然强力减排政策），趋势外推就会严重失准。
\end{warningbox}

\subsection{模型4：多变量回归模型（论文公式21--22）}

\[
\hat{C}_t = \lambda_0 + \sum_{k=1}^{K} \lambda_k X_{k,t}
\]

把GDP、能源消费、煤炭占比等6个变量同时放进回归来预测碳排放。参数通过OLS估计：
\[
\hat{\lambda} = (X^\top X)^{-1} X^\top C
\]

\begin{warningbox}[为什么表现不好？]
18个数据点要估计7个参数（6个变量 + 截距），自由度只有11。变量之间还有相关性（共线性）。结果就是估计出来的系数不稳定，预测误差很大。
\end{warningbox}

\subsection{模型5：随机森林（论文公式23）}

\begin{conceptbox}[随机森林]
一种\textbf{集成学习}方法。核心思想：建很多棵"决策树"，每棵树用不同的数据子集和特征子集来训练，最后把所有树的预测结果\textbf{取平均}。
\end{conceptbox}

\subsubsection{什么是决策树？}

\begin{conceptbox}[决策树]
像一个"20个问题"游戏。每次选一个特征做判断（比如"人均GDP是否超过5万？"），把数据分成两半，然后在每一半继续判断，直到分出最终结果。
\end{conceptbox}

\subsubsection{随机森林为什么叫"随机"？}
两个"随机"：
\begin{enumerate}[leftmargin=2em]
\item \textbf{数据随机}：每棵树用随机抽取的子集训练（有放回抽样，叫\textit{Bagging}）
\item \textbf{特征随机}：每次分裂只考虑随机选出的部分特征
\end{enumerate}

\begin{warningbox}[论文中随机森林为什么崩了？]
随机森林通常需要成千上万个数据点才能发挥优势。论文只有18个数据点，每棵树能分的数据太少，容易"记住"训练数据的噪声而不是学到真实规律。这就是\textbf{过拟合}（下文详述）。RMSE高达5054.08，比朴素模型差3倍。
\end{warningbox}

\subsection{模型6：XGBoost（论文公式24--25）}

\begin{conceptbox}[XGBoost]
\textit{eXtreme Gradient Boosting}。一种\textbf{梯度提升树}方法。与随机森林不同，XGBoost的树是\textbf{一棵一棵按顺序建}的：第一棵树拟合原始数据，第二棵树拟合第一棵树的\textbf{残差}（预测值与真实值的差），第三棵树拟合第二棵树的残差……最终预测是所有树的加和。
\end{conceptbox}

\subsubsection{XGBoost的正则化}

论文公式25中有一个正则项：
\[
\Omega(f_m) = \gamma J + \frac{1}{2}\lambda \|w\|^2
\]
\begin{itemize}[leftmargin=2em]
\item $J$：叶节点数量——节点越多，树越复杂
\item $w$：叶节点的权重值
\item $\gamma$ 和 $\lambda$：正则化系数，控制对复杂度的惩罚力度
\end{itemize}

\begin{tipbox}[正则化的直觉]
就像考试时老师限制你"最多只能用3个公式来解题"。限制越多，你越不容易"死记硬背"特定题目（过拟合），而是被迫学到通用方法。
\end{tipbox}

\begin{warningbox}[XGBoost在论文中的表现]
RMSE=5214.83，MAPE=15.47\%，\textbf{所有模型中最差}（和多变量回归一起垫底）。原因同样是数据太少导致严重过拟合。
\end{warningbox}

\subsection{模型7：集成模型（论文公式26--27）}

\[
\hat{C}_t^{(\text{ens})} = \omega_1 \hat{C}_t^{(\text{AR2})} + \omega_2 \hat{C}_t^{(\text{Trend})} + \omega_3 \hat{C}_t^{(\text{Multi})}
\]

\begin{conceptbox}[集成模型]
把多个模型的预测结果\textbf{加权平均}。权重按照各模型在验证集上的RMSE倒数来确定：
\[
\omega_j = \frac{1/\text{RMSE}_j}{\sum_{j=1}^{3} 1/\text{RMSE}_j}
\]
误差越小的模型权重越大。
\end{conceptbox}

\begin{tipbox}[集成的逻辑]
"三个臭皮匠顶个诸葛亮"。如果不同模型犯了不同的错误，取平均后错误会相互抵消。但如果所有模型犯同样的错误（比如都低估了），集成也没用。
\end{tipbox}

\subsection{七种模型的误差比较总结}

\begin{table}[H]
\centering
\small
\begin{tabularx}{\textwidth}{lcccX}
\toprule
\textbf{模型} & \textbf{RMSE} & \textbf{MAE} & \textbf{MAPE(\%)} & \textbf{一句话点评} \\
\midrule
AR2 & 1323.36 & 1054.50 & 3.32 & \textcolor{green!60!black}{\textbf{最佳}}，简单但稳健 \\
Naive & 1732.09 & 1440.44 & 4.79 & 最简单的基准，居然排第二 \\
Quadratic Trend & 2081.10 & 1805.35 & 5.50 & 趋势外推，尚可接受 \\
Ensemble & 2612.10 & 1772.16 & 5.70 & 集成没带来额外优势 \\
Random Forest & 5054.08 & 4540.33 & 14.27 & \textcolor{red}{\textbf{严重过拟合}} \\
XGBoost & 5214.83 & 5056.28 & 15.47 & \textcolor{red}{\textbf{最差}}，复杂度惩罚 \\
Multivariate OLS & 5959.72 & 2909.50 & 10.13 & 变量太多，自由度不足 \\
\bottomrule
\end{tabularx}
\caption{七种预测模型误差比较}
\end{table}

%% ---------- 12 ----------
\section{评估指标详解}

\subsection{RMSE：均方根误差（论文公式28）}

\[
\text{RMSE} = \sqrt{\frac{1}{n}\sum_{t=1}^{n}(C_t - \hat{C}_t)^2}
\]

\begin{conceptbox}[RMSE]
先算每个预测值与真实值的差，\textbf{平方}后求平均，再开根号。单位与原始数据相同。
\end{conceptbox}

\begin{tipbox}[RMSE的特点]
对\textbf{大误差}惩罚更重（因为先平方了）。如果有个别年份预测偏差特别大，RMSE会明显变大。适合"不能容忍极端错误"的场景。
\end{tipbox}

\subsection{MAE：平均绝对误差（论文公式29）}

\[
\text{MAE} = \frac{1}{n}\sum_{t=1}^{n}|C_t - \hat{C}_t|
\]

\begin{conceptbox}[MAE]
预测误差绝对值的平均。单位与原始数据相同。
\end{conceptbox}

\begin{tipbox}[MAE vs RMSE]
MAE对大误差不那么敏感，更"宽容"。如果有一个年份偏差特别大，RMSE会飙升，但MAE只会温和上升。MAE更\textbf{鲁棒}（robust）。
\end{tipbox}

\subsection{MAPE：平均绝对百分比误差（论文公式30）}

\[
\text{MAPE} = \frac{1}{n}\sum_{t=1}^{n}\left|\frac{C_t - \hat{C}_t}{C_t}\right| \times 100\%
\]

\begin{conceptbox}[MAPE]
用误差占真实值的\textbf{百分比}来衡量，\textbf{无量纲}。3.32\%的MAPE意味着平均预测偏差约为真实值的3.32\%。
\end{conceptbox}

\begin{tipbox}[MAPE的优势]
可以跨不同量纲的数据比较。比如你预测碳排放的MAPE是3\%，预测温度的MAPE是5\%，可以说"碳排放预测得更好"，尽管它们单位不同。
\end{tipbox}

\subsection{MAPE的解读标准}

\begin{table}[H]
\centering
\begin{tabular}{ll}
\toprule
\textbf{MAPE范围} & \textbf{预测质量评价} \\
\midrule
 $< 10\%$ & 高精度预测 \\
 $10\% - 20\%$ & 良好预测 \\
 $20\% - 50\%$ & 合理预测 \\
 $> 50\%$ & 不可接受 \\
\bottomrule
\end{tabular}
\caption{MAPE的一般评价标准}
\end{table}

\begin{examplebox}[论文中MAPE=3.32\%意味着什么？]
AR2模型的MAPE=3.32\%，属于"高精度预测"。而XGBoost的MAPE=15.47\%，虽然按标准还算"良好"，但和AR2比差了将近5倍，在同样数据下这个差距是不可接受的。
\end{examplebox}

%% ---------- 13 ----------
\section{过拟合：论文最核心的方法论教训}

\subsection{什么是过拟合？}

\begin{conceptbox}[过拟合]
模型在\textbf{训练数据}上表现很好，但在\textbf{新数据}上表现很差。本质是模型"死记硬背"了训练数据的噪声和特殊特征，而不是学到了通用规律。
\end{conceptbox}

\begin{examplebox}[一个通俗的类比]
\textbf{欠拟合}：学生太懒，连课本上的例题都不会做。

\textbf{良好拟合}：学生理解了方法，能做课本上的题，也能做考试中的新题。

\textbf{过拟合}：学生把课本上每道题的\textbf{答案}都背下来了，课本上的题全对。但考试时题目稍微一变，就完全不会了。

XGBoost和随机森林在论文中就是第三种情况——在训练数据上可能拟合得很好，但遇到新数据（滚动预测时）就崩了。
\end{examplebox}

\subsection{为什么会过拟合？}

\begin{enumerate}[leftmargin=2em]
\item \textbf{模型太复杂}：参数太多、树太深、特征太多
\item \textbf{数据太少}：模型有足够的能力"记住"每一个数据点
\item \textbf{噪声太多}：数据中的随机波动被模型当成了规律
\end{enumerate}

\begin{warningbox}[论文的极端案例]
18个数据点 + XGBoost（可以建很深的树）= \textbf{教科书级的过拟合}。XGBoost有足够的灵活性把18个点几乎完美地"连"起来，但这些连接方式完全是针对这18个点的，没有任何泛化能力。
\end{warningbox}

\subsection{如何避免过拟合？}

\begin{enumerate}[leftmargin=2em]
\item \textbf{选简单的模型}：像论文中的AR2
\item \textbf{增加数据量}：论文做不到（数据就这么多）
\item \textbf{正则化}：如XGBoost中的 $\gamma$ 和 $\lambda$，但小样本下可能不够
\item \textbf{交叉验证}：但时间序列不太适合标准交叉验证
\item \textbf{用滚动预测评估}：论文的做法，最能反映真实外推能力
\end{enumerate}

\begin{purplebox}[论文的核心方法论启示]
\textbf{"模型适配性优先于模型复杂度"}。不是越复杂的模型越好，而是\textbf{与数据规模匹配}的模型最好。在小样本场景下，简单的线性模型往往完胜复杂机器学习模型。这个结论虽然反直觉，但在学术界已经被反复验证。
\end{purplebox}

%% ============================================================
%%  PART VI
%% ============================================================
\part{可解释性与情景模拟篇}

%% ---------- 14 ----------
\section{SHAP值：让"黑箱"变透明}

\subsection{为什么需要SHAP？}

XGBoost、随机森林等模型被称为"黑箱"——你只知道输入和输出，不知道内部是怎么做决策的。SHAP就是为了解决"黑箱可解释性"问题。

\subsection{什么是SHAP值？}

\begin{conceptbox}[SHAP]
\textit{SHapley Additive exPlanations}。源于博弈论中的Shapley值思想：在一个"合作博弈"中，每个参与者对总收益的贡献是多少？
\end{conceptbox}

\begin{examplebox}[Shapley值的原始场景]
三个人合作做生意赚了100万。A出了资金，B出了技术，C出了人脉。怎么分这100万才公平？Shapley值给出了一种数学上公允的分法：考虑每个人在所有可能的"组合"中的边际贡献。
\end{examplebox}

\subsection{SHAP在机器学习中的应用}

把上面的场景搬到机器学习：
\begin{itemize}[leftmargin=2em]
\item "合作博弈" $\to$ "模型预测"
\item "参与者" $\to$ "特征变量"
\item "总收益" $\to$ "预测值"
\end{itemize}

SHAP值回答的问题是：\textbf{对于某一个预测结果，每个特征贡献了多少？}

论文公式(35)给出了SHAP值的精确定义：
\[
\phi_j(t) = \sum_{S \subseteq F \setminus \{j\}} \frac{|S|!(|F|-|S|-1)!}{|F|!} \left[ \hat{f}_{S \cup \{j\}}(x_{S \cup \{j\}}^t) - \hat{f}_S(x_S^t) \right]
\]

\begin{tipbox}[不用被这个公式吓到]
直觉就是：对于特征 $j$，考虑所有\textbf{不包含} $j$ 的特征子集 $S$，计算"加入 $j$ 后预测值的变化"，然后对所有可能的子集取\textbf{加权平均}。这个加权平均就是特征 $j$ 的SHAP值。
\end{tipbox}

\subsection{SHAP值的三个公理性质}

\begin{enumerate}[leftmargin=2em]
\item \textbf{效率性}：所有特征的SHAP值之和 = 预测值 $-$ 基线值（即"贡献加起来等于总变化"）
\item \textbf{对称性}：如果两个特征对模型的贡献相同，它们的SHAP值也相同
\item \textbf{虚拟性}：如果某个特征对预测没有影响，它的SHAP值为0
\end{enumerate}

\subsection{全局重要性与局部解释}

\begin{itemize}[leftmargin=2em]
\item \textbf{全局重要性}：论文用 $|\phi_j|$ 的均值来衡量。即在所有样本点上，特征 $j$ 的SHAP值绝对值的平均。值越大说明这个特征整体上越重要。
\item \textbf{局部解释}：对于某一次具体预测，看每个特征的SHAP值，知道"这次预测主要被哪些特征驱动"。
\end{itemize}

\subsection{论文中SHAP分析的深刻发现}

\begin{purplebox}[SHAP揭示了XGBoost的"懒惰"]
SHAP显示：人均GDP的重要性\textbf{远超}其他所有变量，城镇化率和技术进步的重要性为\textbf{0}。

\textbf{为什么？}因为在18期数据中，人均GDP和碳排放几乎完美同步增长。XGBoost只需要一个分裂节点（"人均GDP是否超过某阈值？"）就能把低排放年份和高排放年份完美分开。它"懒得"再去挖掘其他变量的作用了。

\textbf{对比STIRPAT}：STIRPAT是多元线性回归，\textbf{强制}所有变量同时入模、同时竞争解释力。正是这种"强制"，让它识别出了煤炭占比的独立弹性。

\textbf{结论}：SHAP分析没有推翻STIRPAT，反而从算法行为的角度证明了——小样本下复杂模型不仅预测差，连机制识别都会"偏废"。
\end{purplebox}

%% ---------- 15 ----------
\section{情景模拟：如果政策不同，未来会怎样？}

\subsection{什么是情景模拟？}

\begin{conceptbox}[情景模拟]
设定不同的\textbf{假设条件}（如不同的减排力度），分别推算未来的可能路径。不是预测"一定会怎样"，而是回答"如果怎样，就会怎样"。
\end{conceptbox}

\begin{examplebox}[类比]
天气预报说"如果明天是晴天，最高温35°C；如果是阴天，最高温28°C"。这不是在预测到底是晴天还是阴天，而是在告诉你不同条件下分别会怎样。情景模拟的逻辑一样。
\end{examplebox}

\subsection{论文的三种情景（论文公式31--34）}

以2022年排放量为起点，设定三种不同的"衰减路径函数" $g^{(s)}(k)$：

\subsubsection{基准情景}
\[
g^{(\text{base})}(k) = 1 - 0.0035k + 0.00015 \cdot \max(k-3, 0)^2
\]
含义：维持现有政策，排放缓慢下降，后期因技术积累降速略快。

\subsubsection{低碳情景}
\[
g^{(\text{low})}(k) = 1 - 0.014k
\]
含义：加速能源结构优化和产业升级，排放以较快\textbf{线性}速率下降。

\subsubsection{深度脱碳情景}
\[
g^{(\text{deep})}(k) = 1 - 0.024k
\]
含义：能源替代与效率提升\textbf{协同强化}，以最快速率压减排放。

其中 $k = t - 2022$ 是距2022年的年数。

\begin{examplebox}[算一算]
到2035年（$k=13$）：
\begin{itemize}[leftmargin=2em]
\item 基准：$g = 1 - 0.0035 \times 13 + 0.00015 \times 100 = 1 - 0.0455 + 0.015 = 0.9695$，排放约为2022年的96.95\%
\item 低碳：$g = 1 - 0.014 \times 13 = 0.818$，排放约为2022年的81.8\%
\item 深度脱碳：$g = 1 - 0.024 \times 13 = 0.688$，排放约为2022年的68.8\%
\end{itemize}
\end{examplebox}

\subsection{情景参数的来源}

论文表4列出了各情景下关键参数的设定：
\begin{itemize}[leftmargin=2em]
\item GDP年均增速：基准约5\%，低碳4.6\%，深度脱碳4.4\%
\item 煤炭占比年均变动：低碳每年降1.1个百分点，深度脱碳降1.6个百分点
\item 能源强度年均下降：低碳2.8\%，深度脱碳3.6\%
\end{itemize}

\begin{tipbox}[参数不是拍脑袋定的]
这些参数参考了《"十四五"现代能源体系规划》和国家统计局近年宏观数据，是有依据的"合理假设"。
\end{tipbox}

\subsection{2035年预测结果}

\begin{table}[H]
\centering
\begin{tabular}{lc}
\toprule
\textbf{情景} & \textbf{2035年碳排放} \\
\midrule
基准 & 39127.33 \\
低碳 & 33013.06 \\
深度脱碳 & 27766.48 \\
\bottomrule
\end{tabular}
\end{table}

\begin{tipbox}[怎么读这个结果？]
从基准到深度脱碳，2035年排放差距超过1.1亿吨。这意味着如果政策力度足够大，可以在2035年前多减排约1.1亿吨——这为政策制定提供了明确的\textbf{量化激励}。
\end{tipbox}

%% ============================================================
%%  PART VII
%% ============================================================
\part{论文写作与方法论篇}

%% ---------- 16 ----------
\section{论文的总体框架设计}

\subsection{五步逻辑链}

论文的结构可以概括为一条清晰的逻辑链：
\begin{enumerate}[leftmargin=2em]
\item \textbf{数据整合}：CEADs + 国家统计局 $\to$ 统一口径的国家级时序数据
\item \textbf{机制识别}：EKC + STIRPAT + 中介分析 $\to$ "为什么排放会变？"
\item \textbf{预测比较}：七种模型滚动预测 $\to$ "哪种模型预测最准？"
\item \textbf{情景模拟}：三种政策情景 $\to$ "不同政策下未来怎样？"
\item \textbf{政策解释}：把统计结果翻译成政策语言
\end{enumerate}

\begin{tipbox}[这个框架好在哪？]
很多论文只做"机制解释"或只做"预测"，这篇把两者\textbf{打通}了——用机制分析的结果来支撑情景设定，用预测比较的结果来选择模拟工具。各章节之间不是孤立的，而是一个整体。
\end{tipbox}

%% ---------- 17 ----------
\section{文献综述的写法}

论文的文献综述按"三条研究路线"组织：
\begin{enumerate}[leftmargin=2em]
\item 碳排放驱动因素研究：恒等式分解 $\to$ 计量回归 $\to$ 系统协同
\item 碳排放预测方法研究：传统统计 $\to$ 机器学习 $\to$ 多模型比较
\item 情景模拟与政策评估研究
\end{enumerate}

最后总结"三点不足"：
\begin{enumerate}[leftmargin=2em]
\item 机制解释与预测精度分离
\item 用样本内拟合评价预测能力（不真实）
\item 情景分析与机制识别脱节
\end{enumerate}

\begin{tipbox}[文献综述的标准写法]
不是简单罗列"张三说了什么、李四说了什么"，而是\textbf{按主题线索}组织，最后指出\textbf{已有研究的不足}，从而自然引出"本文要做什么、为什么有价值"。论文做到了这一点。
\end{tipbox}

%% ---------- 18 ----------
\section{稳健性分析：怎么证明结论"站得住脚"？}

\begin{conceptbox}[稳健性分析]
用不同的方法、不同的假设、不同的数据子集来检验同一个结论。如果结论在各种变化下都成立，就说结论是"稳健的"。
\end{conceptbox}

论文中的稳健性策略包括：
\begin{itemize}[leftmargin=2em]
\item 多模型比较：不只依赖一个模型
\item 滚动预测 vs 样本内拟合：用更严格的评估标准
\item 跨方法互证：STIRPAT和SHAP从不同角度得出一致结论
\item 附录补充：残差诊断、增速对比、滚动相关热力图等
\item 明确区分观测值、预测值、情景值
\end{itemize}

%% ---------- 19 ----------
\section{局限性讨论：诚实的学术态度}

论文明确承认了以下局限：
\begin{enumerate}[leftmargin=2em]
\item 国家级时序样本长度有限（只有18期）
\item 技术进步和治理变量用代理指标，不够精确
\item 聚焦国家总量，未做区域和行业分解
\item 未做严格协整检验（因为样本太小）
\item 情景参数设定仍有一定主观性
\end{enumerate}

\begin{tipbox}[为什么要写局限性？]
在学术论文中，\textbf{主动承认局限}不是减分项，反而是加分项。它说明作者清楚自己研究的边界，没有过度解读结果。审稿人/评委最反感的是"把相关性说成因果"、"把小样本结论当普适真理"。
\end{tipbox}

%% ============================================================
%%  PART VIII
%% ============================================================
\part{附录：关键公式速查表}

\newpage
\section{公式速查表}

\begin{longtable}{p{3cm} p{5.5cm} p{6cm}}
\toprule
\textbf{名称} & \textbf{公式} & \textbf{一句话解释} \\
\midrule
\endfirsthead
\toprule
\textbf{名称} & \textbf{公式} & \textbf{一句话解释} \\
\midrule
\endhead
\bottomrule
\endfoot

碳排放汇总 & $C_t = \sum_{i=1}^{N} C_{i,t}$ & 各省排放加总 \\
人均GDP & $A_t = \text{GDP}_t / P_t$ & 总GDP除以总人口 \\
能源强度 & $EI_t = EC_t / \text{GDP}_t$ & 单位GDP能耗 \\
碳强度 & $CI_t = C_t / \text{GDP}_t$ & 单位GDP碳排放 \\
EKC检验 & $\ln C_t = \alpha_0 + \alpha_1 \ln A_t + \alpha_2(\ln A_t)^2 + \varepsilon_t$ & 倒U型检验 \\
EKC拐点 & $A^* = \exp(-\alpha_1 / 2\alpha_2)$ & 拐点处人均GDP \\
STIRPAT原版 & $\ln I = \ln a + b\ln P + c\ln A + d\ln T + \ln e$ & IPAT的随机化扩展 \\
扩展STIRPAT & $\ln C_t = \beta_0 + \beta_1\ln A_t + \beta_2\ln EI_t + \beta_3 CS_t + \beta_4 UR_t + \beta_5 TP_t + \beta_6 IS_t + \varepsilon_t$ & 六因素模型 \\
置信区间 & $\hat{\beta}_j \pm 1.96 \cdot \text{SE}(\hat{\beta}_j)$ & 95\%置信区间 \\
中介a路径 & $M_t = \gamma_0 + \gamma_1 X_t + u_t$ & X对M的影响 \\
中介c路径 & $Y_t = \delta_0 + \delta_1 X_t + v_t$ & X对Y的总效应 \\
中介c'+b路径 & $Y_t = \theta_0 + \theta_1 X_t + \theta_2 M_t + w_t$ & 直接效应+中介效应 \\
效应分解 & $\delta_1 = \theta_1 + \gamma_1 \cdot \theta_2$ & 总=直接+间接 \\
朴素预测 & $\hat{C}_t = C_{t-1}$ & 随机游走基准 \\
AR2 & $\hat{C}_t = \phi_0 + \phi_1 C_{t-1} + \phi_2 C_{t-2}$ & 二阶自回归 \\
二次趋势 & $\hat{C}_t = \eta_0 + \eta_1 t + \eta_2 t^2$ & 抛物线趋势 \\
多变量OLS & $\hat{C}_t = \lambda_0 + \sum \lambda_k X_{k,t}$ & 多元回归预测 \\
随机森林 & $\hat{C}_t^{(\text{RF})} = \frac{1}{B}\sum_{b=1}^{B} T_b(X_t)$ & 多树取平均 \\
XGBoost & $\hat{C}_t^{(\text{XGB})} = \sum_{m=1}^{M} f_m(X_t)$ & 逐步拟合残差 \\
集成模型 & $\hat{C}_t^{(\text{ens})} = \sum \omega_j \hat{C}_t^{(j)}$ & 加权平均 \\
RMSE & $\sqrt{\frac{1}{n}\sum(C_t - \hat{C}_t)^2}$ & 均方根误差 \\
MAE & $\frac{1}{n}\sum|C_t - \hat{C}_t|$ & 平均绝对误差 \\
MAPE & $\frac{1}{n}\sum\left|\frac{C_t-\hat{C}_t}{C_t}\right|\times100\%$ & 平均百分比误差 \\
情景路径 & $C_t^{(s)} = C_{T_0} \cdot g^{(s)}(t-T_0)$ & 不同政策下的路径 \\
\label{tab:formula}
\end{longtable}

\newpage
\section{核心术语中英对照表}

\begin{longtable}{lll}
\toprule
\textbf{中文} & \textbf{英文} & \textbf{缩写} \\
\midrule
\endfirsthead
\toprule
\textbf{中文} & \textbf{英文} & \textbf{缩写} \\
\midrule
\endhead
\bottomrule
\endfoot

碳达峰 & Carbon Peak & -- \\
碳中和 & Carbon Neutrality & -- \\
碳排放 & Carbon Emissions & -- \\
环境库兹涅茨曲线 & Environmental Kuznets Curve & EKC \\
普通最小二乘法 & Ordinary Least Squares & OLS \\
弹性 & Elasticity & -- \\
半弹性 & Semi-elasticity & -- \\
置信区间 & Confidence Interval & CI \\
显著性水平 & Significance Level & -- \\
自由度 & Degrees of Freedom & df \\
异方差 & Heteroscedasticity & -- \\
多重共线性 & Multicollinearity & -- \\
平稳性 & Stationarity & -- \\
单位根 & Unit Root & -- \\
一阶单整 & Integrated of Order 1 & I(1) \\
伪回归 & Spurious Regression & -- \\
协整 & Cointegration & -- \\
自回归模型 & Autoregressive Model & AR \\
随机森林 & Random Forest & RF \\
梯度提升树 & Gradient Boosting & -- \\
极端梯度提升 & eXtreme Gradient Boosting & XGBoost \\
集成学习 & Ensemble Learning & -- \\
过拟合 & Overfitting & -- \\
正则化 & Regularization & -- \\
滚动预测 & Rolling Forecast & -- \\
情景模拟 & Scenario Simulation & -- \\
SHAP值 & SHapley Additive exPlanations & SHAP \\
对数变换 & Log Transformation & -- \\
标准化 & Standardization & -- \\
中介效应 & Mediation Effect & -- \\
描述性统计 & Descriptive Statistics & -- \\
相关系数 & Correlation Coefficient & -- \\
能源强度 & Energy Intensity & EI \\
煤炭消费占比 & Coal Share & CS \\
城镇化率 & Urbanization Rate & UR \\
第二产业占比 & Secondary Industry Share & IS \\
\label{tab:terms}
\end{longtable}

\vspace{1cm}
\begin{center}
\Large\textcolor{mainblue}{\textbf{—— 全文完 ——}}
\end{center}

\vspace{0.5cm}
\begin{center}
\small\textcolor{gray}{本文档覆盖了原论文中出现的所有核心概念、公式与方法。\\
建议按目录顺序阅读，每学完一个概念后回到论文对应位置对照理解。\\
遇到不理解的地方，多看"例子"和"论文原话解读"两个框。}
\end{center}

\end{document}